% secV2_setting.tex -- Section 2. CONTAINER (27 Oct 2026): text to be written by the author.
% Word budget: 800 words (JPART count: running text, footnotes and citations as printed).
% Outline, numbers and sources: referee_EE_2026-10/ESQUELETO.md, section 2. The archived text
% (_arquivo_texto_IA_2026-10-27/) is not a source.

\section{Institutional and legal setting}\label{sec:institutions}

Childcare for children below age four is not compulsory in Brazil, unlike preschool for four- and five-year-olds, which became compulsory under the 2009 constitutional amendment \citep{brasil2009ec59}. Access to early-childhood education is nevertheless an enforceable right. In 2022, Brazil's Supreme Federal Court held that public childcare and preschool provision can be demanded individually (Theme 548, RE 1,008,166) \citep{stf2022tema548}. Registration consequently captures families' expressed requests for a public place, not the population of all children who might benefit from childcare. The 2014 National Education Plan established a population target of serving at least 50\% of children under four \citep{brasil2014pne}.

S\~ao Paulo's municipal system uses three arrangements: directly operated centres (\emph{rede direta}), municipally owned centres under private management (\emph{rede indireta}), and private centres operating under municipal agreements (\emph{rede conveniada}). Childcare places are divided among Ber\c{c}\'ario I, Ber\c{c}\'ario II, Minigrupo I, and Minigrupo II according to age at the March 31 cut-off. The corresponding adult-to-child ratios are 1:7 for ages 0--11 months, 1:9 for 12--23 months, 1:12 for 24--35 months, and 1:25 for 36--47 months \citep{sme2020in36}. Between 88\% and 89\% of Ber\c{c}\'ario I classes operate exactly at their authorized class size, compared with 23\% to 25\% of Minigrupo II classes (Table~\ref{tab:T6a_modulo}), while the ratio of enrolment to authorized places ranges from 0.93 to 0.996 (Table~\ref{tab:T6b_ratio}). Because assignments must also match a child's age group, vacancies are not necessarily interchangeable across children.

Families can apply throughout the year at no charge. The registration system uses the child's address to identify geographically eligible centres rather than collecting a ranked list of preferred centres. At each centre, applications are ordered chronologically within priority categories, including judicial orders and specified vulnerability or special-needs criteria. Families may also elect to remain registered for a particular centre after being offered a place elsewhere, preserving the original protocol order for that centre under a 2011 municipal ordinance \citep{sme2011portaria3478}. A centre-specific request may therefore remain on the published list despite an alternative offer. Geographical assignment rules have changed over time; we do not impose a fixed distance radius across the sample.

Referrals can occur throughout the year, but the allocation calendar is highly seasonal. Of first referrals dated 2013--2017, 50.5\% occurred in December (Table~\ref{tab:calendar}) and another 15.8\% in January--February, whereas only 5.0\% of first applications occurred in December. We refer to December--February as the main allocation round. A year-end list count therefore reflects the timing of referrals as well as the stock of unresolved applications. Childcare provision also expanded substantially: all-network childcare-stage enrolment increased from 170,720 in 2010 to 365,854 in 2018, across a change in the treatment of Minigrupo II in 2011. Authorized places for children under four rose by 58\% from the first semester of 2013 to the first semester of 2019 (Tables~\ref{tab:T2a_capacidade} and~\ref{tab:R_capcheck}; Figure~\ref{fig:F7}). Meanwhile, annual live births declined by 6.3\% between their 2011 peak and 2018. Between 2013 and 2019, the number of partner-operated and indirectly managed centres increased from 1,223 to 2,004, while directly managed centres declined from 418 to 402 (Table~\ref{tab:T7a_ceis}), and places per eligible child rose from 0.307 to 0.514, with a descriptive decomposition assigning 86\% of this increase to capacity growth and 14\% to the decline in the eligible population (Table~\ref{tab:S6c_cobertura}). The distribution of enrolment recorded in the municipal childcare request data also shifted towards private centres operating under municipal agreements, while the shares of directly managed and indirectly managed municipal centres declined (Table~\ref{tab:T5c_modal}). These trends jointly affect observed enrolment rates and queue counts but do not, by themselves, identify policy effects.

S\~ao Paulo adopted its Municipal Education Plan (PME) under Law 16,271 of September 17, 2015, establishing 13 goals over a ten-year horizon. Goal~5 mandates universal preschool for four- and five-year-olds by 2016 and, during the plan's validity, childcare provision for at least 75\% of children aged 0--3 years and 11 months or 100\% of registered demand, whichever is greater \citep{saopaulo2015pme}. The childcare target exceeds the national plan's population threshold; Law 17,437/2020 later extended the PME goal deadlines by two years \citep{saopaulo2020lei17437}. Strategy 5.1 calls for expansion through direct, indirect, and partnered centres.

The law specifies neither the definition of registered demand, nor a population denominator, nor a rule for reconciling enrolment measures. Subsequent documents introduced monitoring conventions. Ordinance 7,972 of the Municipal Education Secretariat (SME), issued in December 2016, organizes annual monitoring in January--February, with a July round for indicators relying on the School Census, and instructs that indicators be recorded as of December 31 \citep{sme2016portaria7972}. The 2016 baseline report uses December 31, 2015 as its reference date, without a published formula permitting independent reconstruction of Goal~5's compliance measure \citep{sme2016linhabase}. The municipality's 2023 conference review classified Goal~5 as achieved \citep{saopaulo2023conferencia}. Section~\ref{sec:goal} examines the measurement conventions necessary to interpret that classification.
